\section{Conclusão}

Este trabalho comparou a efetividade de suítes PBT equivalentes em duas implementações equivalentes de uma API REST, em Java/Spring Boot e em Node.js/Express, medida pelo \ms{} em \totalExecucoes{} execuções completas de mutação. Quanto à hipótese do estudo, \textbf{$H_0$ foi rejeitada em 11 das 12 comparações}: em todas as seis com os operadores padrão, sempre com o Node.js acima, e em cinco das seis com operadores harmonizados, com o Java acima na suíte fraca e o Node.js acima na suíte forte. \textbf{$H_0$ não foi rejeitada} apenas na suíte forte com 1000 tentativas e operadores harmonizados, em que os dois ecossistemas chegaram a \mshJavaForteMil\%. Com o mutante distinto como unidade, nenhuma diferença harmonizada foi significativa. A diferença entre os ecossistemas existe, portanto, mas depende sobretudo da ferramenta de mutação e da força das propriedades, e não da linguagem.

Os objetivos específicos foram atendidos da seguinte forma:

\begin{enumerate}[label=\textbf{OE\arabic*}, leftmargin=1.6cm, itemsep=1pt]
  \item As 20 regras puderam ser expressas como propriedades equivalentes em jqwik e fast-check, sem adaptação; a suíte fraca verifica 10 delas, e nenhuma de suas propriedades confere o valor do preço.
  \item A força do oráculo foi o fator dominante: com 100 tentativas, o \ms{} subiu de \msJavaFracaCem\% para \msJavaForteCem\% no Java e de \msNodeFracaCem\% para \msNodeForteCem\% no Node.js. Mais tentativas só ajudaram a suíte forte (de 12 a 18 pontos entre 10 e 1000 tentativas, contra menos de 3 na suíte fraca), porque os defeitos de limite só aparecem quando o gerador produz o valor exato e só são reconhecidos quando a propriedade confere esse valor.
  \item Cobertura alta não significou verificação: a suíte fraca cobriu mais de 91\% das linhas do domínio e detectou cerca de metade dos mutantes.
  \item Na suíte fraca, a maioria dos mutantes não detectados está em regras sem propriedade; na forte, em limites exatos que dependem da entrada gerada. Os equivalentes são poucos (\clsJavaForteEquivalente{} no Java e \clsNodeForteEquivalente{} no Node.js).
  \item O custo dependeu da suíte: com a suíte fraca, o StrykerJS foi mais rápido por mutante; com a forte, o PIT foi mais rápido no total e por mutante. A configuração dos dois ecossistemas tem tamanho parecido.
\end{enumerate}

Para quem testa APIs com PBT, a implicação prática é que as propriedades de limite e de neutralidade, comuns na prática \citep{corgozinho2023pbt}, protegem pouco os cálculos e as validações: é preciso escrever ao menos uma propriedade com modelo de referência por regra de cálculo e cobrir as regras de validação, e só então aumentar o número de tentativas. Para quem compara ferramentas ou ecossistemas, a implicação é reportar o \ms{} sobre um conjunto de operadores harmonizado, e não o total de cada ferramenta.

\textbf{Trabalhos futuros.} (i) Replicar o estudo em outras aplicações, para que a unidade amostral seja o programa e não a execução; (ii) aplicar PBT com estado diretamente sobre a camada HTTP, observando as respostas das operações, o que eliminaria a principal causa de oráculo fraco no Java; (iii) estudar estratégias de geração voltadas a limites derivados da especificação como alternativa a aumentar o número de tentativas; (iv) implementar nas duas ferramentas operadores idênticos, inclusive na direção das condições forçadas; e (v) comparar suítes escritas por desenvolvedores sem acesso à especificação completa com as deste trabalho.
